\scanpage{108}
\uline{Dém} L'équivalence de (i) et (ii) est valable, chaque fois qu'on a un fibré vectoriel orthogonal $A$ (rg \unsure{quelc.}) avec forme $d$ non dégénéré, et qu'on regarde \deleted{les} \inserted{un} \deleted{hyp} sous fibré vectoriel de codim $1$ (en supposant $\operatorname{rg}A$ pair \uline{ou} $2$ inv. dans $k$) -- \deleted{quand on \unsure{est}} \inserted{dans} ce dernier cas, inutile de parler de disc. \uline{divisé}, le discriminant ordinaire suffit).

En vertu de $\operatorname{disc}_{\varphi_B}=\operatorname{disc}_{\psi_B}=\inserted{2\operatorname{disc}'_{\det_{\mid B}}}$, \deleted{\unsure{la conv} et} la condition (ii\unsure{)} implique (ii$'$), et la réciproque est vraie si $2$ inv.

\deleted{Comme} Si $D$ est isotrope, il est orthogonal à lui-$=$, donc $D\unsure{\in|\subset}D^\perp=B$. L'inverse est vrai si $2$ est inv. (NB si $2=0$ dans $k$, alors on a \uline{tjrs} $D\subset B$ !); En effet, $D\subset B$ signifie que $D$ est orthogonal à lui-$=$, et donne $2d(x)=\varphi_d(x,x)=0$ Cela implique $2d(x)=0$ ou $d(x)=0$ pour $x\in D$ i.e. $D$ est isotrope.

On sait que les sous-alg. de Borel sont en corr. $1$-$1$ avec les $D$ isotropes \inserted{dans $A^{\unsure{s|\leftrightarrow|\natural}}$}, à \deleted{$D$ cor} une telle $D$
% archive sheet number: 107
